% !TeX root = ../../tesis.tex
% Esquema ilustrativo (decisión del autor 2026-11-04): claves y rutas
% tomadas de VFTModel/src/api/main.py; valores = marcadores de tipo.

\section{Ejemplo de Respuesta: Factor de Desviación \texorpdfstring{($DI$)}{(DI)}}
\label{anx:vft-out-detour}

Esquema ilustrativo de la respuesta del indicador $DI$
(§\ref{subsec:ruta-directa}). Los nombres de las claves y la ruta son los
de \texttt{VFTModel}; los valores son marcadores de tipo. Las cifras de la
corrida 2026-10-04 se reportan en el Capítulo~\ref{ch:analisis-red-actual}.

A diferencia de los demás, este \termino{endpoint} está en la raíz de la \api{},
fuera del prefijo \texttt{/api/v1/network/}.

\begin{lstlisting}[style=estiloPython,
  caption={Esquema de respuesta del endpoint \texttt{/get\_detour\_factor}
           (ilustrativo; \texttt{visualize} = falso).},
  label=lst:out-detour]
{
  "status": "success",
  "parametros": { "muestra_estadistica": <número>,
                  "modo_grafo": "REALISTIC_INTEGRATION",
                  "visualizacion_activa": false },
  "data": [
    { "Origen": "<texto>", "Sistema_Origen": "<sistema>",
      "Destino": "<texto>", "Sistema_Destino": "<sistema>",
      "Sistemas_Involucrados": ["<sistema>"],
      "Ruta_Estaciones": ["<estación>"],
      "Saltos_Topologicos": <número>,
      "Distancia_Recta_km": <número>,
      "Distancia_Red_km": <número>,
      "Factor_Desviacion": <número>,
      "Consideraciones_Reales": {
        "Caminata_1ra_Milla_m": <número>,
        "Caminata_Ultima_Milla_m": <número>,
        "Estacion_Abordaje": "<estación>",
        "Estacion_Descenso": "<estación>" } }
  ]
}
\end{lstlisting}
